% # Espacios compactos

% **Def**: $X$ es compacto si todo cubrimiento por abiertos de $X$ tiene un subcubrimiento finito.


% **Ej**: $\mathbb R$ es compacto, sea:

% $$
% \begin{align*}
% U = \{(n-1,n+1)\}_{n \in \mathbb Z}
% \end{align*}
% $$

% es un cubrimiento de $\mathbb R$ 


% $U$ no tiene subcubrimiento finitos, de hecho $U$ no tiene subcubrimiento propios., muncho menos finitos.


% - b) Si $X$  es un conjunto infinito.
% $(X,\mathcal T_{discreto})$ no es compacto,

% $$
% X = \bigcup_{x \in X} \{x\}
% $$


% $\{\{x\}, x\in X\}$ es un cubrimiento que no tiene subcubrimiento propios,

% -c) $(0,1)$ no es compacto,


% $$
% \begin{align*}
% \bigcup_{(0,1-\frac 1 n): n \in \mathbb N}
% \end{align*}
% $$

% $U$es un vubriemiento abierto de $(0,1)$.  Sean $\{(0,1-\frac 1 {n_1}), (0,1-\frac 1 {n_2}), \ldots, (0,1-\frac 1 {n_k})\}$ un subfamilia finita.


% $$
% \begin{align*}
% \bigcup_{i=1}^k (0,1-\frac 1 {n_i}) = (0,1-\frac 1 {max\{n_1,n_2,\ldots,n_k\}}) \neq (0,1)
% \end{align*}
% $$

% luego $(0,1)$ no es compacto.

% ## compactos

% - a) $(X,\mathcal T$ es compacto , $X$ es finito. 
% - b) $(X,\mathcal T_{\trivial})$ es compacto, $X$ es infinito.

% $(X,\mathcal T)$ $\mathcal T$ finita,  es compacto.


% **Teorema** Heine-Borel: 
% $[0,1]$ es compacto.

% **Dem**: SEa

% $$
% \begin{align*}
% U = \{U_\alpha: \alpha\in I\}
% \end{align*}
% $$ 

% un cubrimiento abierto de $[0,1]$.


% existe $\alpha_0 \in I$ talque $0\in U_{\alpha_0},


% Existe $c_0=\epsilon/2>0$ tq $[0,c_0]\subset U_{\alpha_0}$, luego $[0,c_0]$ esta contenido en la union finita de finitos elemento de $U$.


% SEa 
% $$
% \begin{align*}
% L = \sup\{c \in [0,1]: [0,c] \quad \text{esta contenido en la union de finitos elementos de $U$}\}
% \end{align*}
% $$


% $L$ esta bien definido $L>0$ por que $1\geq L\geq C_0 > 0$ 

% $
% L \in [0,1]$ luego existe $\beta \in I$ talque 

% $$
% \begin{align*}
% L \in U_\beta
% \end{align*}
% $$

% Existe $\epsilon>0$ tq $[L-\epsilon,L] \subseteq U_\beta$. Ahora $L-\epsilon$ no es una cota superor de del conjunto por el cual se conturye $L$ , luego existe $c$ en ese conjunto tq $L-\epsilon < c$. es decir $[0,C]$ esta contenido en la union de finitos elementos de $U$ $U_{\alpha_1},\ldots,U_{\alpha_l}$, Esot es

% $$
% \begin{align*}
% [0,C] \subseteq U_{\alpha_0} \cup U_{\alpha_1} \cup \ldots \cup U_{\alpha_l}
% \end{align*}
% $$

% Luego 

% $$
% \begin{align*}
% [0,L] \subseteq U_{\alpha_0} \cup U_{\alpha_1} \cup \ldots \cup U_{\alpha_l} \cup U_\beta
% \end{align*}
% $$

% Esto muestra que $L$ es en conjunto sobre el se define $L$. Ademas si suponemos que $L<1$ entonces el mismo argumento demuestra que 
% $$
% \begin{align*}
% [0,L+\epsilon] \subseteq  U_\beta
% \end{align*}
% $$
% y por lo tanto $:

% $$
% \begin{align*}
% [0,L+\epsilon]\subseteq \bigcup_{i=1}^l U_{\alpha_i} \cup U_\beta
% \end{align*}
% $$

% entonces $L+\epsilon$ esta en ese conjunto. Lo cual es una contradiccion. Por lo tanto $L=1$  y $1$ esta en ese conjunto. 

% $[0,1]$ es contenido en la union de finitos elementos de $U$.


% **Prop** $X$ compacto y $A$ subconjunto de $X$ cerrado. Entonces $A$ es compacto.

% **Dem**

% Sea $U$ un cubrimiento abierto de $A$. 

% $$
% \begin{align*}
% U = \{U_i \cap A: i \in I \land U_i \subseteq^{ab}  X\}
% \end{align*}
% $$

% $X-A \subseteq^{ab}  X$ es abierto. 

% Consideremao la familia.

% $$
% \begin{align*}
% U' = \{U_i: i \in I\} \cup \{X-A\}
% \end{align*}
% $$

% $U_i$ es un cubrimiento abierto de $X$.


% $$
% \begin{align*}
% \bigcup_{i \in I} U_i \cup A = A
% \end{align*}
% $$
% entonces
% $$
% \begin{align*}
% \bigcup_{i \in I} U_i  \supseteq A
% \end{align*}
% $$

% $$
% \begin{align*}
% X-A \supseteq X-A_i
% \end{align*}
% $$

% entonces

% $$
% \begin{align*}
% \bigcup U' = X
% \end{align*}
% $$

% Como $X$ es compacto existen $U_1, \ldots, U_n, X-A$ tal que 


% $$
% \begin{align*}
% \left(
%     \bigcup_{i=1}^n U_i 
% \right) U (X-A) = X
% \end{align*}
% $$


% entonces intersectando con $A$  tenemos que


% $$
% \begin{align*}
% \bigcup_{i=1}^n U_i \cap A = A
% \end{align*}
% $$

% con $U_1\cap A, \ldots, U_n \cap A$ es una familia finita y por lo tanto $A$ es compacto.


% **Ejemplos**
% - a) $\mathcal C \subseteq^{\text{cerr}} [0,1]$ luego es compacto.
% - $\{\frac 1 n: n \in \mathbb Z^+\}$ no  es compacto con la topologia de subespacio.
% -  $\{\frac 1 n: n \in \mathbb Z^+\} \cup \{0\}$  es compacto con la topologia de subespacio.

% **Prop**

% Si $f:X \to Y$ es continua, y $X$ es compacto, entonces $f(X)$ es compacto.

% **Dem** Sea $U$ un cubrimiento de $f(x)$ por abiertos,asi sea:

% $$
% \begin{align*}
% f^-1(U) = \{f^{-1}(U_i): i \in I\}
% \end{align*}
% $$

% $f^{-1}(U)$ es un cubrimiento abierto de $X$, como $X$ es compacto, existen $U_1,\ldots,U_n$, tal que:

% $$
% \begin{align*}
% X = \bigcup_{i=1}^n f^{-1}(U_i)
% \end{align*}
% $$

% entonces
% $$
% \begin{align*}
% f(X) = \bigcup_{i=1}^n U_i
% \end{align*}
% $$

% Luego $f(X)$ es compacto.

% **Obs** $[a,b]$ es compacto.


% **Prop** $A\subseteq \mathbb R$ cerrado y acotado es compato

% **Dem**

% Como $A$ es acotado, entonces existen $a,b \in \mathbb R$ tq:

% $$
% \begin{align*}
% A\subseteq [a,b]
% \end{align*}
% $$

% Pero $A \subseteq^{cerr} \mathbb R$ entonces $A \cap [a,b] \subseteq^{cerr} [a,b]$ con $A = A \cap [a,b]$. 

% Por la prop anterior (2 veces) $A$ es compacto.

% **Teorema** Si $X$ y $Y$ son compactos, entonces $X \times Y$ es compacto.


% **Dem**


% Sea $\{U_\alpha: \alpha \in I\}$ un cubrimiento abierto de $X \times Y$.


% Si $(x,y) \in X \times Y$ entonces existe $\alpha$ tal que $(x,y) \in U_\alpha$. Ademas existen $U_{(x,y)} \subseteq^{ab}  X$ y $V_{(x,y)} \subseteq^{ab} Y$, $(x,y) \in U_{(x,y)} \times V_{(x,y)}\subseteq U_\alpha$.

% Ahora fijemos $x$ y dejemos variando $y$ 

% $$
% \begin{align*}
% \{U_{(x,y)} \times V_{(x,y)}: y \in Y\} \end{align*}
% $$

% un cubrimiento de $\{x\} \times Y$ quees comapcto, luego existen $y_1,\ldots,y_{n_x}$ tal que:


% $$
% \begin{align*}
% \{x\} \times Y = \left(\bigcup_{i=1}^{n_x} U_{(x,y_i)} \times V_{(x,y_i)}\right) \cap (\{x\} \times Y)
% \end{align*} 
% $$

% entonces:

% $$
% \begin{align*}
% \{x\} \times Y \subseteq \bigcup_{i=1}^{n_x} U_{(x,y_i)} \times V_{(x,y_i)}
% \end{align*}
% $$

% Sea 
% $$
% \begin{align*}
% x \in U_x = \bigcap_{i=1}^{n_x} U_{(x,y_i)} \subseteq^{ab}  X
% \end{align*}
% $$

% Se tiene que 
% $$
% \begin{align*}
% \{x\} \times Y \subseteq U_x \times V_{(x,y_i)} \subseteq \bigcup_{i=1}^{n_x} U_{(x,y_i)} \times V_{(x,y_i)} \subseteq U_\alpha
% \end{align*}
% $$ 

% Notese que $\{U_x, x\in X\}$ es un cubrimiento $X$, como es $X$ es compacto, existen $x_1,\ldots,x_{m}$ tal que:


% $$
% \begin{align*}
% X = \bigcup_{i=1}^m U_{x_i}
% \end{align*}
% $$

% Luego:


% $$
% \begin{align*}
% X \times Y = \bigcup_{i=1}^m U_{x_i} \times Y = \bigcup_{i=1}^m \bigcup_{j=1}^{n_{x_j}} U_{(x_i,y_j)} \times V_{(x_i,y_j)}
% \end{align*}
% $$

% Es decir $X \times Y$ esta contiendoa en uan union finita de $U_{\alpha}$'s, 


% Entonces $X \times Y$ esta contiednos en uan subfamilia finita de $\{U_\alpha\}_{\alpha \in I}$, luego $X \times Y$ es compacto.


% **Teorema** $X \subseteq^{ab}  \mathbb R^n$ es compacto sii $X$ es cerrado y acotado.


% **Dem** ($\impliedby$) E.P.L


% ($\implies$) Veamosque $X$ es acotado. Sea 

% $$
% \begin{align*}
% U = \{B_n(0): n \in \mathbb N\}
% \end{align*}
% $$

% es un cubrimiento de $\mathbb R^n$ Como $X$ es compacto, existen $B_{n_1}(0),\ldots,B_{n_k}(0)$ tal que:

% $$
% \begin{align*}
% X = \bigcup_{i=1}^k B_{n_i}(0)
% \end{align*}
% $$

% Asi $X$ es acotado.


% Falta ver que $X$ es cerrado. Sea $x$ un punto limite de $X$. Suponga que $x \notin X$. Entonce para todo $r>0$ 

% $$
% \begin{align*}
% \overline{B_r(x)} \cap X \neq\emptyset
% \end{align*}
% $$
% Sea 

% $$
% \begin{align*}
% U = \{\mathbb R^n - \overline{B_r(x)}: r>0\}
% \end{align*}
% $$

% donde $\mathbb R^n - \overline{B_r(x)} \subseteq^{ab}  \mathbb R^n$:

% - $U$ cubre a $X$.

% $$
% \begin{align*}
% \bigcup U \supseteq X
% \end{align*}
% $$


% existen $r_1,\ldots,r_l$ tal que:


% $$
% \begin{align*}
% X \subseteq \bigcup_{i=1}^l 
% \mathbb R^n - \overline{B_{r_i}(x)} = \mathbb R^n -  \overline{B_{\min{r_1,\ldots,r_l}}(x)}
% \end{align*}
% $$

% Luego $B_{\min{r_1,\ldots,r_l}}(x)$ es una vecindad de $x$ que no contien puntos de $X$. Es no es posible por que $x$ es un punto limite de $X$. Por lo tanto $X$ es cerrado.
\subsection{Espacios compactos}

\paragraph{\textbf{Definición.}}
$X$ es compacto si todo cubrimiento por abiertos de $X$ tiene un subcubrimiento finito.

\paragraph{\textbf{Ejemplo.}}
$\mathbb R$ \emph{no} es compacto. Sea
\begin{align*}
U=\bigl\{(n-1,n+1)\bigr\}_{n\in\mathbb Z}
\end{align*}
un cubrimiento de $\mathbb R$.  
$U$ no posee subcubrimientos finitos (ni siquiera subcubrimientos propios).

\begin{itemize}
  \item[b)] Si $X$ es un conjunto infinito, entonces $(X,\mathcal T_{\text{discreto}})$ no es compacto porque
    \begin{align*}
    X=\bigcup_{x\in X}\{x\}
    \end{align*}
    es un cubrimiento por abiertos sin subcubrimiento propio (en particular, sin subcubrimiento finito).
  \item[c)] $(0,1)$ no es compacto.  
    El cubrimiento abierto
    \begin{align*}
      \bigl\{(0,1-\tfrac1n)\bigr\}_{n\in\mathbb N}
    \end{align*}
    no posee subcubrimiento finito, pues para cualesquiera $n_1,\dots,n_k$
    \begin{align*}
      \bigcup_{i=1}^{k}\!\bigl(0,1-\tfrac1{n_i}\bigr)=
      \bigl(0,1-\tfrac1{\max\{n_1,\dots,n_k\}}\bigr)\neq(0,1).
    \end{align*}
\end{itemize}

\subsubsection{Compactos básicos}

\begin{itemize}
  \item[a)] Si $(X,\mathcal T)$ es compacto y $X$ es finito, el resultado es trivial.
  \item[b)] El espacio $(X,\mathcal T_{\text{trivial}})$ es compacto aun cuando $X$ sea infinito.
  \item $(X,\mathcal T)$ con $\mathcal T$ finita es compacto.
\end{itemize}

\paragraph{\textbf{Teorema (Heine–Borel).}}
$[0,1]$ es compacto.

\paragraph{\textbf{Demostración.}}
Sea 
\begin{align*}
U=\{U_\alpha\}_{\alpha\in I}
\end{align*}
un cubrimiento abierto de $[0,1]$.  
Existe $\alpha_0\in I$ tal que $0\in U_{\alpha_0}$.  
Hay $\varepsilon>0$ con $[0,\varepsilon]\subseteq U_{\alpha_0}$.  
Por definición de cubrimiento, $[0,\varepsilon]$ está contenido en la unión de finitísimos elementos de $U$.

Sea
\begin{align*}
L=\sup\bigl\{c\in[0,1]\;:\;[0,c]\text{ está contenido en la unión de finitos elementos de }U\bigr\}.
\end{align*}
El número $L$ está bien definido, $0<L\le1$.  
Existe $\beta\in I$ con $L\in U_\beta$; elige $\delta>0$ tal que $[L-\delta,L]\subseteq U_\beta$.  
Entonces $L-\delta$ no es cota superior del conjunto anterior; hay $c>L-\delta$ con $[0,c]$ cubierto por finitos elementos de $U$, digamos $U_{\alpha_1},\dots,U_{\alpha_\ell}$.  
En consecuencia
\begin{align*}
[0,L]\subseteq U_{\alpha_0}\cup U_{\alpha_1}\cup\dots\cup U_{\alpha_\ell}\cup U_\beta.
\end{align*}
Esto contradice la definición de $L$ salvo que $L=1$.  
Por tanto $[0,1]$ admite un subcubrimiento finito y es compacto. \hfill$\square$

\paragraph{\textbf{Proposición.}}
Si $X$ es compacto y $A\subseteq X$ es cerrado, entonces $A$ es compacto.

\paragraph{\textbf{Demostración.}}
Sea $\{V_i\}_{i\in I}$ un cubrimiento abierto de $A$.  
Cada $V_i\subseteq X$ es abierto y $X\setminus A$ es abierto; entonces
\begin{align*}
U'=\bigl\{V_i\bigr\}_{i\in I}\cup\{X\setminus A\}
\end{align*}
es un cubrimiento abierto de $X$.  
Por compacidad de $X$ existe un subcubrimiento finito $V_{i_1},\dots,V_{i_n},X\setminus A$.  
Al intersectar con $A$ se obtiene
\begin{align*}
A=\bigcup_{j=1}^{n}\!\bigl(V_{i_j}\cap A\bigr),
\end{align*}
con lo cual $A$ es compacto. \hfill$\square$

\paragraph{\textbf{Ejemplos.}}
\begin{itemize}
  \item[a)] El conjunto de Cantor $\mathcal C\subseteq[0,1]$ es compacto por ser cerrado en un compacto.
  \item[b)] $\{\tfrac1n\;:\;n\in\mathbb Z^{+}\}$ no es compacto en la topología de subespacio de $\mathbb R$.
  \item[c)] $\{\tfrac1n:n\in\mathbb Z^{+}\}\cup\{0\}$ \emph{sí} es compacto en esa topología.
\end{itemize}

\paragraph{\textbf{Proposición.}}
Si $f\colon X\to Y$ es continua y $X$ es compacto, entonces $f(X)$ es compacto.

\paragraph{\textbf{Demostración.}}
Sea $\{U_i\}_{i\in I}$ un cubrimiento abierto de $f(X)$.  
Entonces $\bigl\{f^{-1}(U_i)\bigr\}_{i\in I}$ cubre $X$.  
Por compacidad existe $U_{i_1},\dots,U_{i_n}$ tales que
\begin{align*}
X=\bigcup_{j=1}^{n}f^{-1}(U_{i_j}),
\end{align*}
luego
\begin{align*}
f(X)=\bigcup_{j=1}^{n}U_{i_j},
\end{align*}
que es subcubrimiento finito. \hfill$\square$

\paragraph{\textbf{Observación.}}
Todo intervalo cerrado y acotado $[a,b]$ es compacto.

\paragraph{\textbf{Proposición.}}
En $\mathbb R$, todo subconjunto cerrado y acotado es compacto.

\paragraph{\textbf{Demostración.}}
Sea $A\subseteq\mathbb R$ cerrado y acotado.  
Existen $a<b$ con $A\subseteq[a,b]$.  
Dado que $A$ es cerrado en $[a,b]$ y $[a,b]$ es compacto, la proposición anterior implica que $A$ es compacto. \hfill$\square$

\paragraph{\textbf{Teorema.}}
Si $X$ y $Y$ son compactos, entonces $X\times Y$ es compacto.

\paragraph{\textbf{Demostración.}}
Sea $\{U_\alpha\}_{\alpha\in I}$ un cubrimiento abierto de $X\times Y$.  
Para cada $(x,y)\in X\times Y$ elige $\alpha$ con $(x,y)\in U_\alpha$ y abiertos $U_{(x,y)}\subseteq X$, $V_{(x,y)}\subseteq Y$ tales que
\begin{align*}
(x,y)\in U_{(x,y)}\times V_{(x,y)}\subseteq U_\alpha.
\end{align*}
Fijado $x\in X$, el conjunto $\{x\}\times Y$ es compacto, por lo que existe una selección finita
\begin{align*}
\{x\}\times Y\subseteq\bigcup_{i=1}^{n_x}U_{(x,y_i)}\times V_{(x,y_i)}.
\end{align*}
Definamos
\begin{align*}
U_x=\bigcap_{i=1}^{n_x}U_{(x,y_i)},
\end{align*}
de modo que $\{x\}\times Y\subseteq U_x\times Y\subseteq\bigcup_{i=1}^{n_x}U_{(x,y_i)}\times V_{(x,y_i)}$.  
La familia $\{U_x\}_{x\in X}$ cubre $X$; por compacidad de $X$ existen $x_1,\dots,x_m$ tales que $X=\bigcup_{j=1}^{m}U_{x_j}$.  
Por consiguiente
\begin{align*}
X\times Y
&=\bigcup_{j=1}^{m}U_{x_j}\times Y
 =\bigcup_{j=1}^{m}\bigcup_{i=1}^{n_{x_j}}U_{(x_j,y_i)}\times V_{(x_j,y_i)},
\end{align*}
una unión finita de elementos de $\{U_\alpha\}$.  
Así, $X\times Y$ es compacto. \hfill$\square$

\paragraph{\textbf{Teorema (Heine–Borel en $\mathbb R^n$).}}
Un subconjunto $X\subseteq\mathbb R^n$ es compacto si y sólo si es cerrado y acotado.

\paragraph{\textbf{Demostración.}}
\begin{itemize}
  \item[$(\Rightarrow)$] Si $X$ es compacto, el cubrimiento por bolas $U=\{B_k(0)\}_{k\in\mathbb N}$ admite subcubrimiento finito, por lo que $X$ está contenido en una bola cerrada; luego es acotado.  
    Para cerrar, sea $x$ punto límite de $X$ y supón \(x\notin X\).  Los abiertos
    \begin{align*}
      U_r=\mathbb R^n\setminus\overline{B_r(x)},\qquad r>0,
    \end{align*}
    cubren $X$.  De un subcubrimiento finito se obtendría una bola $\overline{B_{\rho}(x)}$ sin puntos de $X$, contradiciendo que $x$ es límite.  Así, $X$ es cerrado.
  \item[$(\Leftarrow)$] Si $X$ es cerrado y acotado, se encierra en un rectángulo producto de intervalos cerrados, que es compacto por el teorema anterior; como subespacio cerrado de un compacto, $X$ es compacto.
\end{itemize}
\hfill$\square$